\section{从闭合电动势到一段导线的端电压}

电磁感应中的电路计算要同时描述两件事：外场怎样推动电荷，以及电荷重新分布后形成怎样的端电压。闭合回路的总电动势由磁通变化确定，但这一结果不能直接分配给任意一段导线。特别是在变化磁场中，把整个回路的电动势都画在某个电阻旁边，虽然有时能算出总电流，却可能算错电压表读数。


在准静态条件下，将电场分成电荷产生的保守部分 $\bm E_s=-\nabla V$ 与外磁场变化产生的部分 $\bm E_f$。这种分解须结合已给定的外场与边界约定使用。导线本身以速度 $\bm u$ 运动时，沿从 $a$ 到 $b$ 的实际导线定义

\[\varepsilon_{ab}=\int_a^b(\bm E_f+\bm u\times\bm B)\cdot\dd\bm l.\]
若该段电流的正向也是 $a\to b$，局部欧姆定律沿线积分给出
\[
 \int_a^b\bm E_s\cdot\dd\bm l+\varepsilon_{ab}=i_{ab}R_{ab},
 \qquad
 \int_a^b\bm E_s\cdot\dd\bm l=-\int_{V_a}^{V_b}\dd V=V_a-V_b.
\]
将保守场积分代入并移项，得到
\begin{equation}V_a-V_b=i_{ab}R_{ab}-\varepsilon_{ab}.\label{em2:branch}\end{equation}
第一项是电阻降，第二项是推动正电荷的非静电作用。在开路的运动棒中，电流为零，端电压与动生电动势互相平衡；在有电流的弧形导线中，两项都不能省去。反向走同一条路径时，电流、电动势和电压的有向量同时反号，电阻本身不变。


式\eqref{em2:branch}也说明为什么必须先标方向。若沿 $a\to b$ 算得 $i<0$，应保留这个负号，最后再说明实际电流由 $b$ 流向 $a$。过早把电动势、电流都取绝对值，会破坏节点电流方程。两个端点相同的支路应具有相同的 $V_a-V_b$，但支路电动势可以不同，正是这一差别驱动环流。


\fig{em2:branchpic}{沿 $a\to b$，非静电作用给出电势升高，电阻给出电势下降}{\begin{tikzpicture}

\fill(0,0)circle(2pt)node[left]{$a$};\fill(6,0)circle(2pt)node[right]{$b$};\draw[world](0,0)--(1,0) (2,0)--(3,0) (4,0)--(6,0);\draw[world](1,-.25)rectangle(2,.25);\node[above]at(1.5,.25){$R_{ab}$};\draw[world](3.5,0)circle(.5);\draw[flow](3.15,0)--(3.85,0);\node[below]at(3.5,-.5){$\varepsilon_{ab}$};\draw[flow](.3,.9)--(2.3,.9)node[midway,above]{$i_{ab}$};\node at(3,-1.5){$V_b=V_a+\varepsilon_{ab}-i_{ab}R_{ab}$};

\end{tikzpicture}}
图\ref{em2:branchpic}把电流和分布电动势都按$a\to b$规定正向。沿箭头经过电阻产生电阻降，而非静电作用向正电荷供能；两者之差才对应同一端点电势差。换向时应同时改变电流、电动势及端电势差的符号。



\subsection{导线内部的电势分布与支路功率}

端点关系是沿导线积分得到的，若题目询问中间点的电势，必须保留局部关系。用弧长 $s$ 表示位置，单位长度电阻为 $r(s)$，沿导线的非静电推动作用为 $e(s)=(\bm E_f+\bm u\times\bm B)\cdot\bm t$，则

\[\frac{\dd V}{\dd s}=e(s)-i\,r(s).\]
电流在没有分支的细导线中处处相同，局部电阻降却可随材料、截面改变。电势的极大值、极小值出现在右端为零且两侧导数变号之处。因此即使两端电势相等，内部也可能存在显著的电势起伏；“导体等电势”只适用于静电平衡，不能直接用于有持续非静电驱动和电流的长导线。


以绕轴匀速转动的均匀径向导棒为例，取 $s$ 自转轴向外、$\bm B\parallel\bm\omega$。图\ref{em2:local-potential}中 $e(s)=B\omega s$，单位长度电阻为 $r_0$，因而

\[
V(s)-V(0)=\int_0^s(B\omega\sigma-ir_0)\dd\sigma
=\left[\frac12B\omega\sigma^2-ir_0\sigma\right]_0^s
=\frac12B\omega s^2-ir_0s.
\]
开路时第二项为零；接负载后要减去沿程电阻降。若 $ir_0=B\omega L/2$，则
\[
 \frac{\dd V}{\dd s}=B\omega(s-L/2),\qquad
 \frac{\dd^2V}{\dd s^2}=B\omega>0.
\]
两端恰等电势，内部却在 $s=L/2$ 达到最低电势；代入上式，
$V(L/2)-V(0)=B\omega L^2/8-B\omega L^2/4=-B\omega L^2/8$。图中的曲线由局部推动与耗散共同决定，不能把整根运动导体当成等电势体。

\fig{em2:local-potential}{同一转棒在开路与有负载时的电势曲线；端点等电势不代表内部等电势}{\begin{tikzpicture}

\begin{scope}
\draw[world](0,0)--(3,0);\fill(0,0)circle(2pt)node[below]{$O$};\fill(3,0)circle(2pt)node[below]{$P$};\draw[flow](1,.4)--(2.3,.4)node[midway,above]{$i,\ s$};\draw[flow](3,.2)--(3,1.2)node[right]{$\bm u$};\draw[->,red,thick](1.5,0)--(2.5,0)node[midway,below]{$\bm u\times\bm B$};\node at(1.1,1.5){$\odot B,\omega$};\node at(1.5,-1){$e(s)=B\omega s$};
\end{scope}
\begin{scope}[xshift=5.5cm]
\draw[axis](-.2,0)--(4.7,0)node[right]{$s/L$};\draw[axis](0,-1.2)--(0,2.5)node[above]{$[V(s)-V(0)]/(B\omega L^2)$};\draw[blue,thick,domain=0:1,samples=60]plot({4*\x},{4*.5*(\x)^2});\draw[red,thick,domain=0:1,samples=60]plot({4*\x},{4*(.5*(\x)^2-.5*\x)});\draw[guide](2,0)--(2,-.5);\node[below]at(2,-.5){$1/2$};\node[below]at(4,0){$1$};\node[blue]at(3.2,2.3){开路};\node[red]at(3.5,-1){两端短接};\node[left]at(0,-.5){$-1/8$};\node[left]at(0,2){$1/2$};
\end{scope}

\end{tikzpicture}}

对支路关系乘电流，得到 $i(V_a-V_b)+i\varepsilon_{ab}=i^2R_{ab}$。左侧第一项是端口电功率，第二项是分布驱动向电流提供的功率；它可以来自运动部件的机械功，也可以来自维持变化磁场的外部装置。若第二项为负，该支路的感应作用正在吸收电能。对全网所有支路求和，节点电势的功率项因电流守恒相消，剩下所有驱动功率之和等于发热率。这是复杂感应网络的整体校验。


开路时 $i=0$，导线内部的静电场沿线抵消非静电推动，因而 $V(s)-V(0)=\int_0^s e(\sigma)\dd\sigma$。接上负载以后，静电场不再完全抵消，电势曲线要减去累积电阻降。若只是改变电阻率而运动与磁场不变，电动势分布不变，但电流和电势分布都可变化。这个区分可帮助分析多接点转棒、非均匀磁场中的导棒和局部电压测量。



\section{空间电场积分与电路等效}

轴对称变化磁场可以直接提供支路的感应积分。圆柱磁区半径为 $a$，内部 $\dot B=k$、外部磁场可忽略时，按正法线取向有
\[E_\theta(r)=-\frac{kr}{2}\quad(r<a),\qquad E_\theta(r)=-\frac{ka^2}{2r}\quad(r>a).\]
沿同心圆弧积分，用 $\dd l=r\dd\theta$；沿径向积分则为零。沿一条偏离圆心的弦积分，必须先把环向电场投影到弦方向，不能把弧长公式套用到直线。路径跨越磁区边界时要分段，边界处电场连续，积分结果也应连续。


将闭合电路分成若干弧段，每一段都可用电阻与有向电动势串联表示。这个等效只保留指定端点与指定路径上的关系，不表示感生电场真的集中在一个理想电源中。尤其是要增加新的引线或比较另一对端点时，应回到空间积分，不能任意挪动等效电源。


对于静止的电阻网络，选定独立回路后，将式\eqref{em2:branch}沿回路相加，电势差望远镜式相消，便得到电阻降之和等于总感应电动势。节点处则满足流入与流出电流之和相等。电流的建立要求回路闭合；开路电容虽然不允许稳态传导电流通过，仍能通过短暂充电建立端电压，须结合电荷守恒处理。


在均匀静磁场中刚性平移的导线，动生积分只与端点有关：
\[\varepsilon_{ab}=(\bm u\times\bm B)\cdot(\bm r_b-\bm r_a).\]
同一刚体内两条连接相同端点的路径，其平移动生电动势相同。可是一条在运动、另一条固定的支路并不满足这个前提；滑动接触问题必须逐段使用真实速度。把静止接触板也赋予移动圆弧的速度，是常见的建模错误。


\fig{em2:field-chord}{环向电场在直线路径上的投影；磁区外仍存在感生电场}{
\begin{tikzpicture}[scale=1.05]
\fill[blue!5](0,0)circle(1.45);
\draw[world](0,0)circle(1.45);
\draw[guide](0,0)circle(2.05);
\draw[axis](-2.8,0)--(2.8,0)node[right]{$x$};
\draw[flow,red](-2.5,.65)--(2.5,.65);
\fill(-2.5,.65)circle(1.5pt)node[above]{$P$};
\fill(2.5,.65)circle(1.5pt)node[above]{$Q$};
\fill(-1.296,.65)circle(1.5pt)node[above left]{$M$};
\fill(1.296,.65)circle(1.5pt)node[above right]{$N$};
\draw[dim](-.25,0)--(-.25,.65)node[midway,left]{$h$};
\draw[guide](0,0)--(.9,.65)node[midway,below]{$r$};
\fill(.9,.65)circle(1.5pt);
\draw[flow,blue](.9,.65)--(1.42,-.07)node[right]{$\bm E_f$};
\draw[flow](0,0)--(-1,-1.05)node[midway,left]{$a$};
\node[align=center]at(.35,-.7){$\odot\ \dot B$\\$k>0$};
\node[below]at(0,-2.1){$PM$、$NQ$在磁区外，$MN$在磁区内};
\end{tikzpicture}}

图\ref{em2:field-chord}中取直线路径$P\to Q$，其高度为$h>0$，微元为$\dd\bm l=\dd x\,\bm e_x$。按逆时针定义$\bm e_\theta$，有$r^2=x^2+h^2$和$\bm e_\theta\cdot\bm e_x=-h/r$，因而
\[
 \bm E_f\cdot\dd\bm l=-E_\theta(r)\frac{h}{r}\,\dd x.
\]
这一步投影把环向电场变为沿弦的推动作用，不能以圆弧长代替$\dd x$。由图中交点$M,N$把积分分成三段，内段代入$E_\theta=-kr/2$，两外段代入$E_\theta=-ka^2/(2r)$。在$r=a$处两式同为$-ka/2$，所以分段不会引入电场跳变。固定端点的测量支路若改走别的路径，也应重新按路径投影积分。



\section{有向支路关系怎样组成完整网络}

设网络有若干节点，将每个节点标一个电势，任意选择一处为零电势。对每条支路指定正电流方向，使用 $i_{ab}=(V_a-V_b+\varepsilon_{ab})/R_{ab}$。随后在每个独立节点写电流守恒，所有方程中的 $V_a-V_b$ 都由同一组节点电势给出，因此不会出现同一路径被重复计算的矛盾。改变零电势只给全部电势加同一个常量，不改变电流或任何可测量结果。


两条并联支路都取 $a\to b$ 为正，电阻与电动势分别为 $R_1,\varepsilon_1$ 和 $R_2,\varepsilon_2$。若外部没有取流，$i_1+i_2=0$，解得

\[V_a-V_b=-\frac{\varepsilon_1/R_1+\varepsilon_2/R_2}{1/R_1+1/R_2},\qquad i_1=\frac{\varepsilon_1-\varepsilon_2}{R_1+R_2}.\]
这里两支路电动势必须使用相同的端点方向。若把第二条支路电动势按反向计算却仍代入这个式子，会把差错误地变成和。两支路电动势相同、外部不开路时，网络也不一定没有电流；相同电动势只能说明它们之间没有由差值驱动的内部环流。


\fig{em2:parallel-branches}{两支路采用相同端点正向；内部环流来自电动势之差}{\begin{tikzpicture}

\draw[world](0,-1.1)--(0,1.1) (6,-1.1)--(6,1.1);\fill(0,0)circle(2pt)node[left]{$a$};\fill(6,0)circle(2pt)node[right]{$b$};
\foreach \yy/\jj in {1.1/1,-1.1/2}{\draw[world](0,\yy)--(1,\yy) (2,\yy)--(3,\yy) (4,\yy)--(6,\yy);\draw[world](1,{\yy-.2})rectangle(2,{\yy+.2});\node[above]at(1.5,{\yy+.2}){$R_\jj$};\draw[world](3.5,\yy)circle(.5);\draw[flow](3.15,\yy)--(3.85,\yy);\node[above]at(3.5,{\yy+.5}){$\varepsilon_\jj$};}
\draw[->,blue,thick](4.3,1.1)--(5.3,1.1)node[midway,above]{$i_1$};\draw[->,red,thick](5.3,-1.1)--(4.3,-1.1)node[midway,below]{$-i_2=i_1$};\node at(3,-2.4){图示 $\varepsilon_1>\varepsilon_2$，两支路正向仍均取 $a\to b$};

\end{tikzpicture}}
图\ref{em2:parallel-branches}中，第二支路的实际电流反向，但其参考电动势仍按 $a\to b$ 积分。沿上路前进、下路返回，总电动势为 $\varepsilon_1-\varepsilon_2$，串联总电阻为 $R_1+R_2$，与节点法所得 $i_1$ 完全一致。



对多回路网络，可选若干独立回路电流求解，但每条公共支路的真实电流是相邻回路电流的有向差。计算发热时先求这个真实电流再平方，不能把两个回路电流的平方分别相加。若要检验解，节点电流守恒与每个独立回路的电压关系都要成立，只满足其中一个回路并不足够。


\fig{em2:shared-branch}{两个相邻回路在公共电阻中反向经过，公共支路真实电流是有向差}{\begin{tikzpicture}

\draw[world](0,0)rectangle(6,2.6);\draw[world](3,0)--(3,.9) (3,1.7)--(3,2.6);\draw[world](2.8,.9)rectangle(3.2,1.7);\node[right]at(3.25,1.3){$R_c$};
\draw[->,blue,thick](1.4,2)arc(90:-220:.6);\node[blue]at(1.4,1.3){$I_1$};\draw[->,red,thick](4.6,2)arc(90:-220:.6);\node[red]at(4.6,1.3){$I_2$};\draw[flow](3.5,2.3)--(3.5,1.8)node[right]{$i_c$};\node at(3,-.6){$i_c=I_1-I_2$};\node at(3,-1.2){$P_c=R_c(I_1-I_2)^2$};

\end{tikzpicture}}
图\ref{em2:shared-branch}的两个回路均取顺时针，上述公共支路取向下为正。右回路在此向上经过，因此实际电流是 $I_1-I_2$，发热中的交叉项不可丢弃。这里的回路电流只是求解变量，直接进入欧姆定律和热量的始终是每段导线中的真实电流。



网络等效时，关闭独立感应驱动用于求端口电阻，只是在线性方程中取驱动为零，不表示真的撤去磁场后机械状态也必然保持原来。含移动导体的端口电阻若与机械反馈同时变化，需要另列运动方程；本讲有限内阻电压表的计算假定导线位置与速度均由外界给定，不包含表负载反作用改变运动的效果。



\subsection{多支路端口等效与零电流支路}

三个以上支路接到同一对节点时，不宜先猜每个小回路电流。各支路均从 $A$ 指向 $B$，令 $U=V_A-V_B$，则 $i_j=(U+\varepsilon_j)/R_j$。若没有外接电流，节点方程为 $\sum_j i_j=0$，于是

\[U=-\frac{\sum_j\varepsilon_j/R_j}{\sum_j1/R_j}.\]
电势差是各支路电动势的电导加权结果，并非算术平均。求出 $U$ 后再还原各支路电流；零电流支路满足 $U+\varepsilon_j=0$，它的电阻降为零，但端电势差一般不为零。把这条支路临时断开不会改变网络状态，正是电桥平衡和零示测量能够简化计算的原因。


\fig{em2:zero-branch}{三支路同电阻的例子：中支路零电流而端电势差非零}{\begin{tikzpicture}

\draw[world](0,2.5)--(6,2.5) (0,0)--(6,0);\node[left]at(0,2.5){$A$};\node[left]at(0,0){$B$};
\foreach \xx/\lab in {1/0,3/{\varepsilon_0},5/{2\varepsilon_0}}{\draw[world](\xx,2.5)--(\xx,1.8) (\xx,1)--(\xx,.7) (\xx,.1)--(\xx,0);\draw[world]({\xx-.2},1)rectangle({\xx+.2},1.8);\node[left]at({\xx-.25},1.4){$R$};\draw[world](\xx,.4)circle(.3);\draw[flow](\xx,.6)--(\xx,.2);\node[below]at(\xx,-.2){$\lab$};}
\draw[->,blue,thick](.5,.6)--(.5,2)node[midway,left]{$\varepsilon_0/R$};\draw[->,blue,thick](5.5,2)--(5.5,.6)node[midway,right]{$\varepsilon_0/R$};\node at(3,3.2){$U=V_A-V_B=-\varepsilon_0$};\node at(3,-1.1){$i_2=(U+\varepsilon_0)/R=0$};

\end{tikzpicture}}
图\ref{em2:zero-branch}将三支路取同一电阻，电动势按向下依次为 $0,\varepsilon_0,2\varepsilon_0$。电导加权给出 $U=-\varepsilon_0$，中支路电流为零，而左右支路仍有等大反向环流。断开中支路不改变另外两支路的状态；“零电流”约束的是支路的总推动 $U+\varepsilon_2$，不能据此把节点电势差置零。



若外接负载电流为 $I$，节点条件改为 $\sum_j i_j+I=0$。整理后，端口表现为一个开路电压与等效内阻 $r_p=(\sum_j1/R_j)^{-1}$。这个等效只回答端口电压、电流与功率；求某条内部导线中间点电势、某一支路热量时，仍须返回内部网络。等效负载不变并不保证内部各电源或感应驱动承担的功率不变。


采用磁通法确定回路电动势时，可先选一个独立回路组。每个闭合回路的正向应与相应曲面法线配对，公共导线在两个相邻回路中常按相反方向经过。不同回路的外磁通不一定相等，也不一定具有相同变化率。先把这些几何关系变成带符号的驱动，再列网络方程，可以避免将邻接磁区的面积重叠计算。


对于磁通测量，电容的最终电压和零电流表的平衡条件都提供有用信息。但只有在充电已稳定、开关没有改变孤立节点总电荷等条件满足时，才可直接用静态端口关系。过程中的电容电流仍需要电路方程；本讲先把端口和几何确定清楚，充电时间的系统计算在暂态电路中展开。


对给定的电阻支路网络，端口方程可直接写为
\[
U=U_{\rm oc}-r_p I,\qquad
U_{\rm oc}=-\frac{\sum_j\varepsilon_j/R_j}{\sum_j1/R_j},\qquad
r_p=\left(\sum_j\frac1{R_j}\right)^{-1}.
\]
这里 $I$ 从 $A$ 流入外负载。接入 $R_L$ 后再联立 $U=IR_L$，便得到下一图的串联等效；求内部热量仍须返回各支路求 $i_j$，端口输出不能代表内部总耗散。



\fig{fv:f2:port}{规定负载路径后，将端口等效为开路电压与串联内阻}{\begin{tikzpicture}

\draw[world](0,1.25)--(0,2.3)--(1,2.3) (2.2,2.3)--(4,2.3);\draw[world](1,2.05)rectangle(2.2,2.55);\node[above]at(1.6,2.55){$R_{\rm th}$};\draw[world](0,.85)circle(.4);\draw[world](0,.45)--(0,0)--(4,0);\node at(0,1.05){$+$};\node at(0,.65){$-$};\node[left]at(-.5,.85){$U_{\rm oc}$};
\fill(4,2.3)circle(2pt)node[above]{$A$};\fill(4,0)circle(2pt)node[below]{$B$};\draw[dashed](4,2.3)--(5,2.3)--(5,1.6) (5,.7)--(5,0)--(4,0);\draw[world](4.75,.7)rectangle(5.25,1.6);\node[right]at(5.3,1.15){$R_L$};\draw[flow](3,2.8)--(4.3,2.8)node[midway,above]{$I$};

\end{tikzpicture}}
图\ref{fv:f2:port}用开路电压 $U_{\rm oc}=V_A-V_B$与串联电阻表示线性网络对指定端口的作用，负载电流的正向从 $A$流入负载、由 $B$返回。接入 $R_L$后，端口电流为 $I=U_{\rm oc}/(R_{\rm th}+R_L)$，端电压为 $IR_L$；开路时虽然电流为零，端电压一般不为零。在含感生电场的装置中，$U_{\rm oc}$还必须对应已规定的引线路径。若改动测量回路使包围磁通改变，就已改变端口描述，不能照搬原等效电源。



\section{电压表与引线构成一个测量回路}

电压表不是脱离电路的抽象数字。它的两根引线、表内电阻及被测对象共同构成闭合路径。若外磁场随时间变化，两条连接相同端点的引线路径可能包围不同磁通，因而得到不同读数。这不违背电势单值性：不同的是测量回路中非保守电场的积分，而非同一个静电势差出现了两个数值。


\fig{em2:lead-paths}{同一对端点的两条测量路径；两表均按从$a$到$b$计有向读数}{
\begin{tikzpicture}[scale=.85]
\fill[blue!6](-2.5,-1.3)rectangle(2.5,1.3);
\draw[world](-2.5,0)--(-2.5,1.3)--(-.35,1.3)
 (.35,1.3)--(2.5,1.3)--(2.5,0);
\draw[world](-2.5,0)--(-2.5,-1.3)--(-.35,-1.3)
 (.35,-1.3)--(2.5,-1.3)--(2.5,0);
\draw[world](0,1.3)circle(.35)node{$V_1$};
\draw[world](0,-1.3)circle(.35)node{$V_2$};
\fill(-2.5,0)circle(2pt)node[left]{$a$};
\fill(2.5,0)circle(2pt)node[right]{$b$};
\draw[flow,blue](-1.9,1.3)--(-1,1.3)node[midway,above]{$g_1$};
\draw[flow,red](-1.9,-1.3)--(-1,-1.3)node[midway,below]{$g_2$};
\node at(0,.1){$\otimes\ \bm n,\ \Phi_{12}$};
\node at(0,-2.05){$\Gamma=g_1+(-g_2)$：沿上路去、下路回};
\end{tikzpicture}}

图\ref{em2:lead-paths}中的两条路径都由$a$指向$b$，因此两支路的端电势项相同。设局部表体的感应可忽略，表内电阻降为$M_j=i_jR_{Vj}$，便有
\[
 M_1=V_a-V_b+\varepsilon_{g_1},\qquad
 M_2=V_a-V_b+\varepsilon_{g_2}.
\]
先作两式之差，再沿图中的闭合方向$\Gamma$使用法拉第定律，得到固定引线时
\[
 M_1-M_2=\int_{g_1}\bm E_f\cdot\dd\bm l
          -\int_{g_2}\bm E_f\cdot\dd\bm l
        =-\frac{\dd\Phi_{12}}{\dd t}.
\]
上路从左向右、下路反向返回，组成顺时针环路，故图中正法线指向纸里。读数差的符号须与这一法线一起确定。比较两次接线时还假定测量分流可忽略，或两次接线没有改变端点电势；若负载效应显著，必须分别重算被测网络。把两根引线贴近，正是缩小图中两路之间的面积。


设表支路由 $a$ 经引线和表体走向 $b$，表体电阻为 $R_V$，其余引线电阻忽略。把该支路的分布电动势记为 $\varepsilon_g$，则

\[i_VR_V=V_a-V_b+\varepsilon_g.\]
若表体局部尺寸很小、表内感应效应可以忽略，仪表按电流标出的读数就是左端 $i_VR_V$，而不是自动等于 $V_a-V_b$。只有引线的感应积分可以忽略时，通常的电压测量公式才直接适用。两根引线相互贴近或绞合，是减小所围面积、抑制这类附加读数的办法。


电阻很大的电压表可以近似不改变被测电路电流，但不能因此删掉引线的感应电动势。反过来，表内电阻有限时，表支路会分流，必须把它作为真正支路联立求解。计算顺序是先画出全部连接路径，写出各支路电动势，再使用节点电流守恒；最后根据表体位置确定读数。若引线也在运动，积分中还应包含它的 $\bm u\times\bm B$，即使外磁场恒定也可能影响读数。



\subsection{测量路径的方向与误差控制}

把两条不同引线路径 $g_1,g_2$ 都由被测正端走向负端，它们的积分之差等于沿 $g_1$ 前进、再沿 $g_2$ 返回的闭合环流。固定引线的差值为 $-\dd\Phi_{12}/\dd t$，其中 $\Phi_{12}$ 的正向由这个闭合行进方向决定。这样可以直接比较读数差，而不必先求空间每一点的静电势。若两次测量对被测电路的分流不可忽略，则还要计入两次不同负载对节点电势的改变。


例如两对紧贴在一起的引线包围窄条面积 $A_g$，磁场近似均匀变化，附加电动势量级为 $A_g|\dot B|$。若容许误差为 $\delta U$，应使 $A_g|\dot B|\le\delta U$，即 $A_g\le\delta U/|\dot B|$；这给出引线布置的面积要求，仅增大表内阻不能替代它。在运动磁场装置附近，引线既可能扫过磁通，也可能被感生电场驱动，两种效应可以在同一个闭合测量回路中统一核算。


端电势差的计算还需要空间场的对称条件。知道整个圆周的总环流，不等于知道任意弦上的电场积分；轴对称时才可断定圆周同半径处环向场大小一致。正方形四边均分总电动势，则需要外场、几何与所选感生场边界条件具有相应旋转对称性，仅仅四条边等长还不够。题目如果没有这类条件，应直接对测量闭合回路用法拉第定律，避免无依据地把电动势平均分配。


静电势与仪表读数的区别不会改变基尔霍夫定律的物理内容。对真实支路把感应积分写全以后，节点电势仍能消去，回路方程仍然成立。问题通常出在把一段分布电动势遗漏了，或者把只适用于某个端口的等效电路用来回答另一条空间路径的问题。



\section{电容、断接操作与孤立导体的电荷}

电容器两极板的电荷为 $\pm q$，约定正极板到负极板的电压为 $u=q/C$。在充电已完成的支路中，电流可以为零而电压不为零。若磁通仍以恒定速率变化，回路的感应电动势可以由电容电压平衡；磁场停止变化以后，这个平衡条件随之改变，不能沿用停止前的电动势。


取走一根跨接导体或者改变开关时，应先识别哪些导体此后彼此绝缘。每个孤立导体连同接在它上面的所有极板，其总电荷守恒；单独某块极板上的电荷却可以经连接导线转移。两个电容并联后的共同电压，来自孤立节点的总电荷除以电容之和，而不是来自两个初始电压的简单平均。


真实导线有很小电阻，电荷重新分配会耗散能量。即使把电阻最后取向零，有限的再分配热量也不必趋于零，因此不能以“电阻忽略”为由同时强加全过程电场能守恒。理想模型若还包含电感，电压与电流可能振荡而不自动稳定；题干说“最终稳定”时，需要明确存在微小耗散。后面的暂态电路将把这些短暂过程展开求解。


电路图只表达连接关系，几何图表达路径与磁通。处理电磁感应测量题，需要两者相互核对：几何图决定每段电动势的符号和大小，电路图决定电流与电荷如何分配。结果还应接受三个检查：反转磁场变化率是否使有向读数反号；电阻趋于相同是否恢复对称性；测量支路开路或短路时是否回到合理极限。



\subsection{重新连接后的节点电荷与能量}

\fig{em2:charge-nodes}{电容重新连接前后：虚线框标出应守恒总电荷的上导体}{
\begin{tikzpicture}[scale=.82]
\begin{scope}
\node at(1.5,3){连接前};
\draw[world](0,2.3)--(1,2.3) (2,2.3)--(3,2.3);
\draw[world](0,2.3)--(0,1.35) (0,1.05)--(0,0)
 (3,2.3)--(3,1.35) (3,1.05)--(3,0) (0,0)--(3,0);
\draw[world](-.35,1.35)--(.35,1.35) (-.35,1.05)--(.35,1.05)
 (2.65,1.35)--(3.35,1.35) (2.65,1.05)--(3.35,1.05);
\draw[world](1,2.3)--(1.8,2.65);
\node at(1.5,1.85){$S$};
\node[left]at(-.35,1.2){$C_1$};
\node[right]at(3.35,1.2){$C_2$};
\node[above]at(.1,2.3){$q_{10}$};
\node[above]at(2.9,2.3){$q_{20}$};
\end{scope}
\draw[flow](4,1.3)--(4.8,1.3);
\begin{scope}[xshift=6cm]
\node at(1.5,3.45){闭合后与外界绝缘};
\draw[blue,dashed,rounded corners](-.5,1.2)rectangle(3.5,2.7);
\draw[world](0,2.3)--(3,2.3)
 (0,2.3)--(0,1.35) (0,1.05)--(0,0)
 (3,2.3)--(3,1.35) (3,1.05)--(3,0) (0,0)--(3,0);
\draw[world](-.35,1.35)--(.35,1.35) (-.35,1.05)--(.35,1.05)
 (2.65,1.35)--(3.35,1.35) (2.65,1.05)--(3.35,1.05);
\node at(1.5,2.95){$N:\ Q=q_{10}+q_{20}$};
\node[left]at(-.35,1.2){$C_1$};
\node[right]at(3.35,1.2){$C_2$};
\node[below]at(1.5,0){$M$};
\end{scope}
\end{tikzpicture}}

图\ref{em2:charge-nodes}中，闭合$S$以后两个上极板与连线组成同一个孤立导体$N$。电荷可以在两块上极板之间转移，但不能越过电容绝缘层离开$N$。因此应守恒的是框内总电荷，而非各块极板原有电荷。以下用$U=V_N-V_M$计算共同电压，初始电荷均按各上极板的有向电荷代入；闭合过程保留微小电阻以保证最终稳定。


设两电容的上极板在连接前分别带电 $q_{10},q_{20}$，下极板总电荷相反。将上极板相连、下极板相连且不接外部电源，微小电阻使最终稳定。上导体总电荷 $Q=q_{10}+q_{20}$，所以共同电压 $U=Q/(C_1+C_2)$，每块上极板的最终电荷为 $C_jU$。若初始极性相反，必须先给电荷带符号，再相加，不能把两个电荷绝对值机械相加。


图\ref{em2:charge-nodes}虚线框的下边界从两极板之间穿过，恰将两块上极板及连线包括在内。最终电压相同意味着 $q_1/C_1=q_2/C_2=U$，配合 $q_1+q_2=Q$ 才解出共同电压。因此两个初始电压通常应按电容加权：

\[
U=\frac{C_1U_1+C_2U_2}{C_1+C_2},\qquad
q_1=\frac{C_1(C_1U_1+C_2U_2)}{C_1+C_2}.
\]
若 $U_1=-U_2$，只有再满足 $C_1=C_2$ 时共同电压才为零；初始电压的绝对值相等并不足够。


用电容初始电压 $U_1,U_2$ 表示，上述再分配热量可由能量差化简为
\[Q_J=\frac12C_1U_1^2+\frac12C_2U_2^2-\frac12(C_1+C_2)U^2
=\frac{C_1C_2}{2(C_1+C_2)}(U_1-U_2)^2.\]
它非负，且只有两初始电压相同时为零。若变化磁场在重新分配期间仍持续做功，应重新列含感应驱动的方程，不能再把初末电场能差全部当成热。本讲取走导体题明确在随后停止磁场变化，正是为了把最终孤立导体电荷约束与持续供能区分清楚。


计算完成后，可以把结果代回某条与求解时不同的路径。弧形电路中用一段求出端电势，再沿另一段计算同一电势差，应一致；测量支路则把求出的节点电势加上引线感应积分，还原表内电阻降。这种检查能够区分代数错误和路径定义错误，比只看最终答案量纲更有针对性。


